% benchmarks/llp_forward/_shared/notes/leptonic_ground_truth.tex
% Part of HEPBench (GPL v3, see LICENSE). GROUND-TRUTH design/physics note --
% NOT for the agent sandbox. A worked SOLUTION KEY for the forward-LLP reach
% problem: it teaches, from the Lagrangian to the exclusion contour, how this
% class of problem is solved, one pipeline step at a time. Organised around the
% sequential steps of the analysis (theory input -> flux -> production ->
% propagation -> assembly -> results), so a reader can learn the method from the
% ground up -- and so the step boundaries expose what tool scaffolding this class
% of problem wants (see the scaffolding section).
\documentclass[11pt]{article}
\usepackage[margin=1in]{geometry}
\usepackage{amsmath,amssymb,mathtools}
\usepackage{graphicx}
\usepackage{booktabs}
\usepackage{listings}
\usepackage{xcolor}
\usepackage[colorlinks=true,allcolors=blue]{hyperref}
\usepackage{cleveref}
\usepackage{slashed}
\lstset{basicstyle=\footnotesize\ttfamily,columns=fullflexible,
        backgroundcolor=\color{gray!6},frame=lines,breaklines=true,
        xleftmargin=1em}
\newcommand{\mphi}{m_\phi}
\newcommand{\Ma}{\mathcal{M}}
\title{\textbf{Forward long-lived-particle reach: a worked solution for the
muon-philic scalar}\\[3pt]
\large From the Lagrangian to the exclusion contour, one pipeline step at a time
(\texttt{llp\_forward})}
\author{HEPBench authors}
\date{\today}
\begin{document}
\maketitle

\begin{abstract}
\noindent This note is a \emph{solution key}. It works the forward long-lived-particle
(LLP) reach problem for a muon-philic scalar $\phi$ end to end --- a scalar
produced in the decays of forward-going mesons at the LHC and decaying, after a
macroscopic flight, inside a downstream detector --- and it is organised the way
the analysis actually runs: a short \emph{theory input} (the decays, the width,
the lifetime), then the four sequential \emph{steps} of the pipeline
(\cref{sec:pythia}~forward parent flux $\to$ \cref{sec:production}~production of
$\phi$ $\to$ \cref{sec:prop}~propagation and acceptance $\to$
\cref{sec:assembly}~assembly and contour), then the \emph{results} and what they
teach. Each step is computed from first principles and validated (the decay
widths against explicit spinor/loop-integral evaluations; the production matrix
elements by Lorentz invariance and a Ward identity to $10^{-32}$). Two physics
lessons fall out and recur throughout: the $(\mphi,g)$ plane splits at the dimuon
threshold $2m_\mu$ into a below-threshold \emph{diphoton lobe} and an
above-threshold \emph{kaon wedge}, separated by a $\sim\!10^3\times$ width cliff;
and the dominant \emph{producer} climbs $K\!\to\!D_s\!\to\!J/\psi$ with mass ---
a real production staircase --- yet at LHC-forward exposure only the kaon step
clears $N_{\rm sig}=3$, so higher-mass reach is \emph{rate}-limited, not
channel-limited. A closing section reads the solution back into
\emph{tool scaffolding}: what each pipeline step asks of a toolchain
(\cref{sec:scaffolding}). All values marked \texttt{[PIN]} are synthetic
placeholders, deliberately matching no real experiment (anti-recall).
\end{abstract}

\tableofcontents

% =====================================================================
\section{The problem, and how this note solves it}
\label{sec:problem}

We want the region of the $(\mphi,g)$ plane that a forward detector can exclude
for a muon-philic scalar,
\begin{equation}
\mathcal{L} \supset -\tfrac12 \mphi^2\phi^2 - g\,\phi\,\bar\mu\mu ,
\label{eq:lag}
\end{equation}
a real scalar coupling to muons with a \emph{scalar} (not pseudoscalar) Yukawa
$g$. One coupling controls both production and decay, and both are exactly
$\propto g^2$ with $g$-independent kinematics, so a single Monte-Carlo sample
reweights across the entire $g$ axis (\cref{sec:assembly}). The observable is a
displaced $\phi$ decay --- $\phi\to\mu^+\mu^-$ above the dimuon threshold, or
$\phi\to\gamma\gamma$ below it.

\paragraph{The master formula, read as a pipeline.} The reach is a product of
factors, and each factor is one step of this note:
\begin{equation}
N_{\rm sig}(\mphi,g)
= N_{\rm int}\;
\underbrace{\Big[\textstyle\sum_M \Phi_M \otimes \tfrac{d\Gamma_{M}}{dx}\Big]}
  _{\substack{\text{production}\\ \text{(\cref{sec:pythia}\ flux,\ \cref{sec:production}\ decay)}}}\;
\underbrace{P_{\rm dec}(c\tau)\,\mathcal{A}}
  _{\substack{\text{propagation}\\(\cref{sec:prop})}}\;
\underbrace{\mathrm{BR}_{\rm vis}(\mphi)}
  _{\substack{\text{theory input}\\(\cref{sec:decay})}}\, ,
\label{eq:master}
\end{equation}
with $N_{\rm int}$ the number of primary $pp$ interactions, $\Phi_M$ the forward
parent-meson flux, $d\Gamma_M/dx$ the exact emission spectrum,
$P_{\rm dec}$ the in-volume decay probability, $\mathcal{A}$ the geometric and
visible-final-state acceptance, and $\mathrm{BR}_{\rm vis}$ the branching to the
search channel. The plan is simply to compute and validate each factor, then
assemble \cref{eq:master} and scan it. A reader can follow the same order to
learn the method; a tool-builder can read the same boundaries as an interface map
(\cref{sec:scaffolding}).

% =====================================================================
\section{Theory input: the two decays, the width, and the lifetime}
\label{sec:decay}

Before any Monte Carlo, the model must be solved as a decay problem: what are the
partial widths, the total width (which sets the lifetime), and the visible
branching fractions? The scalar has \emph{two} visible channels, and both are
mandatory --- one tree, one loop.

\subsection{The dimuon width (tree level)}
\label{sec:decay:mumu}

For the scalar coupling of \cref{eq:lag} the tree-level partial width is
\begin{equation}
\boxed{\;\Gamma(\phi\to\mu^+\mu^-)=\frac{g^2\,\mphi}{8\pi}
\left(1-\frac{4m_\mu^2}{\mphi^2}\right)^{3/2}
=\frac{g^2\,\mphi}{8\pi}\,\beta^3\;}
\label{eq:width}
\end{equation}
with $\beta=\sqrt{1-4m_\mu^2/\mphi^2}$, open for $\mphi>2m_\mu=0.211$~GeV. The
$\beta^3$ (not $\beta^1$!) is the $p$-wave threshold factor of a scalar current
$\bar\mu\mu$; a pseudoscalar current $\bar\mu i\gamma^5\mu$ would give $\beta^1$.
Getting this power right is the single most common way to be wrong by a factor
near threshold. \Cref{eq:width} is checked in \texttt{production.py::validate}
against an explicit two-body Dirac-spinor computation to $<10^{-10}$.

\subsection{The diphoton width (one loop) --- and the only channel below threshold}
\label{sec:decay:gg}

The same $g\,\phi\bar\mu\mu$ vertex generates $\phi\to\gamma\gamma$ through a
muon loop, present at \emph{all} masses and, below the dimuon threshold, the
\emph{only} decay. Omitting it is simply wrong. The width is the Higgs-portal
form with the Yukawa $m_\mu/v\to g$,
\begin{equation}
\Gamma(\phi\to\gamma\gamma)=\frac{\alpha^2 g^2\,\mphi^3}{256\,\pi^3 m_\mu^2}
\,\big|A_{1/2}(\tau)\big|^2,\qquad
A_{1/2}(\tau)=\frac{2[\tau+(\tau-1)f(\tau)]}{\tau^2},\quad
\tau=\frac{\mphi^2}{4m_\mu^2},
\label{eq:gammagg}
\end{equation}
with $f(\tau)=\arcsin^2\!\sqrt\tau$ below threshold, and the exact limits
$A_{1/2}\to4/3$ as $\tau\to0$ and $A_{1/2}(1)=2$. \Cref{eq:gammagg} is checked
against an independent 2-D Feynman-parameter integral
$A_{1/2}=4\!\int_0^1\!dx\!\int_0^{1-x}\!dy\,(1-4xy)/(1-4xy\tau)$ to
$1.5\times10^{-6}$ (\texttt{production.py::gamma\_gammagamma}). This is the one
place in the whole solution where an \emph{exact} result requires a loop; the
tree symbolic tools cannot supply it, so it is derived by hand
(\cref{sec:scaffolding}).

\subsection{Total width, visible branching, lifetime --- and the threshold cliff}
\label{sec:decay:total}

The lifetime is set by the \emph{total} width. Both channels carry exactly one
$g\,\phi\bar\mu\mu$ vertex, so both are $\propto g^2$ and
\begin{equation}
\Gamma_{\rm tot}=g^2\big[\Gamma_{\rm ref,\mu\mu}(\mphi)+\Gamma_{\rm ref,\gamma\gamma}(\mphi)\big],
\qquad
c\tau(g)=\frac{\hbar c}{g^2\,\Gamma_{\rm ref}^{\rm tot}(\mphi)},
\qquad \lambda=\beta\gamma\,c\tau ,
\label{eq:ctau}
\end{equation}
with $\Gamma_{\rm ref}$ the $g^2$-stripped width and $\hbar c=1.9732687\times
10^{-16}$~GeV\,m. The visible branching to a given search,
$\mathrm{BR}_{\rm vis}=\Gamma_{\rm vis}/\Gamma_{\rm tot}$, is
\emph{$g$-independent} (the $g^2$ cancels) --- which, together with the
$g$-independent kinematics, is what will let one sample reweight the whole $g$
axis (\cref{sec:assembly}).

\paragraph{Two regimes, split at $2m_\mu$.} \Cref{fig:lobe} (right) makes the
partition concrete. Above threshold the loop is a $<\!5\times10^{-4}$ correction
($\mathrm{BR}_{\mu\mu}=0.9995$ already at $0.215$~GeV), so the $\mu\mu$ picture is
unchanged within tolerance. Below threshold $\mathrm{BR}_{\gamma\gamma}=1$ and
the width is loop-suppressed to $\Gamma_{\rm ref}\sim10^{-10}$--$10^{-8}$~GeV, a
very long lifetime. The clean signature of the total-width treatment is the
\textbf{threshold cliff}: crossing $2m_\mu$ the width jumps $\sim\!10^3\times$
(loop $\to$ tree), the lifetime collapses, and (as we will see in the contour,
\cref{sec:results}) the below-threshold $\gamma\gamma$ lobe hands the plane to the
above-threshold $\mu\mu$ wedge within a few MeV. This one cliff is the reason the
final reach has two disconnected pieces.

% =====================================================================
\section{Step 1 --- the forward parent flux from Pythia}
\label{sec:pythia}

The scalar is not produced directly; it is radiated in the decays of ordinary
forward mesons. Step~1 is therefore to generate the \emph{parent} flux $\Phi_M$:
the lab momenta and per-collision rates of the mesons that can emit a $\phi$. We
use \textsc{Pythia}~8, configured as in Hostert \emph{et al.}
(arXiv:2507.15271) with taus replaced by our muon parents. Below is the exact
card, explained block by block.

\begin{lstlisting}
Tune:pp = 14                     ! (a) Monash 2013 tune
Beams:idA = 2212                 ! (b) p p beams
Beams:idB = 2212
Beams:frameType = 1              ! (b) collider frame
Beams:eCM = 13600.               ! (b) sqrt(s) = 13.6 TeV        [PIN]
SoftQCD:all = on                 ! (c) soft QCD = forward production
SoftQCD:nonDiffractive = on
SoftQCD:singleDiffractive = on
SoftQCD:doubleDiffractive = on
HardQCD:all = off                ! (c) hard QCD OFF for the forward sample
Onia:all = on                    ! (c) quarkonia on (completeness)
321:mayDecay = off               ! (d) K+-  stable: harvest at production
411:mayDecay = off               ! (d) D+-  stable
431:mayDecay = off               ! (d) Ds+- stable
Print:quiet = on                 ! (e) quiet / bookkeeping
Next:numberCount = 0
Init:showChangedSettings = off
Init:showChangedParticleData = off
Check:event = off
\end{lstlisting}

\paragraph{(a) Tune.} \texttt{Tune:pp = 14} is Monash~2013, the standard
minimum-bias/underlying-event tune; it fixes the hadronization and
multiparton-interaction parameters that set the light-meson and charm yields.

\paragraph{(b) Beams.} Protons in the collider frame at $\sqrt{s}=13.6$~TeV (LHC
Run~3). The beam energy is a \texttt{[PIN]}; a beam-dump sibling would use
\texttt{frameType = 2} (fixed target) at lower $E_{\rm beam}$.

\paragraph{(c) Which QCD --- the key physics choice.} Forward production
(small angle $=$ small $p_T$) is dominated by \emph{soft} QCD, so
\texttt{SoftQCD:all} is on and \texttt{HardQCD} off. This matters: the hard
$2\to2$ $c\bar c$ process produces \emph{central}, high-$p_T$ charm outside the
forward acceptance --- we verified a hard-$c\bar c$ sample with
\texttt{pTHatMin}${=}1$~GeV yields essentially zero forward parents. Forward
charm is \emph{low}-$p_T$ charm, which soft QCD produces at a small rate.
\texttt{Onia:all} is on for completeness (quarkonia are a negligible forward
source here; the charmonium normalization is handled in a separate sample).

\paragraph{(d) Stable parents --- the harvesting trick.} We set $K^\pm,D^\pm,
D_s^\pm$ (and, in the full run, the vectors) \texttt{mayDecay = off}. Being
stable they appear as final-state particles and are harvested directly at their
\emph{production} four-momenta; we then apply the exact $M\to\mu\nu\phi$ decay
ourselves (\cref{sec:production}). This is cleaner than a forced-decay-then-
reweight scheme, and it makes the flux \emph{reusable}: the branching $\kappa_M$
and the detector acceptance are both applied downstream.

\paragraph{(e) Normalization.} After generation we read the generated cross
section $\sigma_{\rm gen}$ (\texttt{infoPython().sigmaGen()}). Each harvested
parent carries a per-inelastic-collision weight
\begin{equation}
w_M = \frac{\sigma_{\rm gen}}{N_{\rm gen}\,\sigma_{\rm inel}},\qquad
\sigma_{\rm inel} = 80~\text{mb}~\text{[PIN]},
\label{eq:wnorm}
\end{equation}
so $\sum w_M$ is the mean number of that parent per $pp$ collision. For the soft
sample $\sigma_{\rm gen}\simeq\sigma_{\rm inel}$.

\paragraph{Forward acceptance folded downstream.} At generation we keep parents
within a \emph{generous} cone $\theta<50$~mrad of either beam (several times the
detector half-angle); the true acceptance is applied later, on the \emph{scalar}
and its decay products (\cref{sec:prop}), because the $\phi$ is emitted at an
angle to the parent. Folding acceptance downstream, not at generation, is exactly
what makes one flux reusable across geometries --- the property the whole
geometry scan (\cref{sec:results}) exploits.

\paragraph{The yield, measured.} A $10^6$-event soft run (793~s) gives, per
collision, the full parent set (the light vectors are harvested in the same run;
$J/\psi,\psi(2S)$ come from the separate onia sample):
\begin{center}\footnotesize
\begin{tabular}{lcccccccc}
\toprule
 & $K$ & $\rho^0$ & $\omega$ & $\phi$ & $D$ & $D_s$ & $J/\psi$ & $\psi(2S)$\\
\midrule
forward/coll & 3.56 & 2.95 & 2.88 & 0.147 & 0.034 & 0.011 & $2.4\times10^{-3}$ & $2.2\times10^{-4}$\\
\bottomrule
\end{tabular}
\end{center}
Kaons and light vectors are abundant and robustly modeled; forward charm is
$\sim\!100\times$ rarer, so it drives the statistics budget and carries the
well-known forward-charm modeling uncertainty (low-$x$ PDFs, intrinsic charm;
cf.\ arXiv:2507.15271). Since our run \emph{is} the reference by construction,
this is a documented systematic, not a correctness issue. \Cref{fig:flux} shows
the resulting spectra. \emph{(Implementation:
\texttt{\_shared/generate\_parent\_flux.py}.)}

\begin{figure}[t]
\centering
\includegraphics[width=\textwidth]{fig_flux_spectra.pdf}
\caption{Forward parent-meson lab energy (left) and angle (right) spectra per
$pp$ collision, from the $10^6$-event Pythia run. Kaons dominate; charm is
$\sim$100$\times$ rarer and slightly more forward-peaked in energy.}
\label{fig:flux}
\end{figure}

% =====================================================================
\section{Step 2 --- production: from parent meson to scalar}
\label{sec:production}

Step~2 turns each harvested parent into a weighted $\phi$. It has two exact
matrix elements (pseudoscalar and vector parents), a kinematic reduction that
makes those matrix elements collapse to \emph{one energy spectrum per mass}, and
two data products (a rate and a shape) that feed the rest of the pipeline.

\subsection{The pseudoscalar amplitude $M\to\mu\nu\phi$}
\label{sec:me}

The scalar is radiated off the muon leg of the ordinary leptonic decay
$M\to\mu\nu$. Because $\phi$ couples \emph{only} to muons, there is a
\emph{single} tree diagram: the meson annihilates through the weak current into
$\mu\nu$, and $\phi$ attaches to the internal (off-shell) muon line. (Contrast
the pseudoscalar-ALP case, where gauge invariance forces a second diagram off
the neutrino leg; the pure Yukawa scalar has no such subtlety.) With
$p_1=p_\mu,\,p_2=p_\nu,\,p_3=p_\phi$ and the off-shell muon propagator
denominator $D=(p_1+p_3)^2-m_\mu^2=2\,p_1\!\cdot p_3+\mphi^2$, the spin-summed
squared amplitude is $\overline{|\Ma|^2}=T/D^2$ (at $g=C_M=1$), with $T$ the
standard trace. The differential rate in $x=2E_\phi^*/m_M$ (the scalar energy
fraction in the $M$ rest frame) and $c=\cos\theta^*$ follows from PDG three-body
phase space,
\begin{equation}
\frac{d^2\Gamma}{dx\,dc}
= C_M^2\;\frac{\overline{|\Ma|^2}\,J(x)}{(2\pi)^3\,32\,m_M^3},
\qquad C_M = \frac{G_F}{\sqrt2}|V_{M}|f_M ,
\label{eq:me}
\end{equation}
with $J$ the Dalitz Jacobian and $f_M,|V_M|$ the decay constant and CKM factor.
The branching coefficient at $g=1$ is
\begin{equation}
\kappa_M(\mphi) = \mathrm{Br}(M\to\mu\nu\phi)\big|_{g=1}
= \frac{C_M^2\,\Gamma_3(\mphi)}{\Gamma_{\rm tot,M}},\qquad
\Gamma_3=\!\int\! \frac{d^2\Gamma}{dx\,dc}\,dx\,dc .
\label{eq:kappa}
\end{equation}
The trace closed form is checked against an explicit Dirac-spinor evaluation at
random Dalitz points ($<10^{-10}$), the $(x,c)$ quadrature against an independent
Dalitz-plane integrator ($<2\times10^{-3}$), and Lorentz invariance under random
boosts ($<10^{-9}$), all in \texttt{production.py::validate}. The three open
parents:
\begin{center}\footnotesize
\begin{tabular}{lccccl}
\toprule
$M$ & $m_M$ [GeV] & $f_M$ [GeV] & $|V|$ & $\tau_M$ [s] & window $\mphi<m_M-m_\mu$\\
\midrule
$K^\pm$  & 0.493677 & 0.1557 & $|V_{us}|{=}0.2243$ & $1.238\times10^{-8}$ & $0.388$ \\
$D^\pm$  & 1.86966  & 0.2120 & $|V_{cd}|{=}0.2210$ & $1.033\times10^{-12}$ & $1.764$ \\
$D_s^\pm$& 1.96835  & 0.2499 & $|V_{cs}|{=}0.9750$ & $5.040\times10^{-13}$ & $1.863$ \\
\bottomrule
\end{tabular}
\end{center}
The pion is absent: $m_\pi-m_\mu=0.034$~GeV $<2m_\mu$, so a $\phi$ from a pion is
too light to decay to $\mu\mu$. \emph{(Implementation:
\texttt{\_shared/production.py}, \texttt{\_shared/parents.py}.)}

\subsection{The vector amplitude $V\to\mu^+\mu^-\phi$ (two diagrams)}
\label{sec:vector}

The scalar is radiated off the muon legs not only of pseudoscalar leptonic decays
but of \emph{vector}-meson dimuon decays,
$V(P)\to\mu^+(p_1)\,\mu^-(p_2)\,\phi(p_3)$, for $V\in\{\rho^0,\omega,\phi(1020),
J/\psi,\psi(2S)\}$, where $V\to\gamma^*\to\mu\mu$ (VMD) and $\phi$ attaches to
\emph{either} muon leg --- \emph{two} interfering diagrams, against the single
diagram of the pseudoscalar case. The only hadronic input is the dimensionless
$\mathrm{BR}(V\to\mu\mu)$: the $V$--$\gamma^*$ coupling cancels in the ratio, so
\begin{equation}
\kappa_V(\mphi)=\mathrm{BR}(V\to\mu\mu)\times R_V(\mphi),\qquad
R_V(\mphi)=\frac{\Gamma(V\to\mu\mu\phi)}{\Gamma(V\to\mu\mu)}\bigg|_{g=1},
\label{eq:kappaV}
\end{equation}
exactly parallel to $\kappa_M$ and computed in the \emph{same} exact
matrix-element framework --- no form factors. The conserved current gives a Ward
identity ($\epsilon\to P$ vanishes), the primary cross-check: it holds to
$3\times10^{-32}$, and the 2-body normalization matches the closed form to
$10^{-16}$ (\texttt{production\_vector.py::validate}).

\subsection{The kinematic reduction: why a single energy spectrum, not the full
matrix element, is what the flux needs}
\label{sec:reduction}

It is worth pausing on \emph{why} the production collapses to one tabulated curve
per mass, because it is the design decision that shapes the whole pipeline.
Picture the decaying parent at rest: a spin-0 kaon is a structureless blob with
no internal axis, so its decay cannot prefer any direction --- whatever comes out
is sprayed uniformly over $4\pi$. All the non-trivial physics is in \emph{how the
released energy is shared} among the products. Now ask what a detector that sees
\emph{only} the scalar can possibly learn from a given decay: exactly two numbers
--- the $\phi$'s energy (which fixes its boost and lab decay length
$\lambda=\beta\gamma\,c\tau$) and its flight direction (which fixes whether it
enters the volume). The partner $\mu$ and $\nu$ are never reconstructed at the
production point, so any detail of \emph{how they} share the recoil is invisible
to the reach. That is the whole reason one curve suffices.

Make it precise. The rest-frame decay has two internal (Dalitz) variables: the
scalar energy fraction $x=2E_\phi^*/m_M$ and a muon helicity cosine
$c=\cos\theta^*$. The direction is a separate, trivial degree of freedom (the
overall orientation of a spinless decay). So the density factorises,
\begin{equation}
\frac{d\Gamma}{dE_\phi^*\,d\Omega^*}
=\underbrace{\frac{d\Gamma}{dx}}
  _{\substack{\text{energy: the one}\\\text{non-trivial curve}}}
\;\times\;
\underbrace{\frac{1}{4\pi}}_{\substack{\text{direction:}\\\text{isotropic}}},
\qquad
\frac{d\Gamma}{dx}\equiv\int_{-1}^{1}\frac{d^2\Gamma}{dx\,dc}\,dc ,
\label{eq:factorize}
\end{equation}
and this is an \emph{identity}, not a modelling step, for two reasons that map
onto the intuition above. \textbf{(i)}~The direction is isotropic and independent
of $x$ because a spin-0 unpolarised parent decays rotationally invariantly ---
$\overline{|\Ma|^2}$ depends only on the internal Dalitz energies, never on where
the decay points. \textbf{(ii)}~The variable $c$ we integrate away only sets how
$\mu$ and $\nu$ arrange themselves \emph{at fixed $\phi$ energy}; it never enters
$\phi$'s momentum, so $d\Gamma/dx$ keeps every degree of freedom the flux can
see. One therefore samples $x\sim d\Gamma/dx$ by inverse-CDF, throws an isotropic
direction, and boosts to the lab (\cref{sec:assembly}), recovering the exact
scalar momentum distribution with \emph{no} run-time matrix-element evaluation
--- the phase-space integral is done once, offline, to build the curve.

\paragraph{The vector parents: the same reduction, one honest caveat.} The two
interfering diagrams of \cref{sec:vector} change the \emph{shape} of the curve
$d\Gamma/dx$, computed in the same exact framework, but not the fact that a curve
suffices. The real subtlety is that the parent now carries spin: a $V$ \emph{can}
be polarised, and a polarised $V$ would correlate the $\phi$ direction with its
polarisation axis --- breaking the isotropy that reason (i) gave us for free. So
for vectors the energy-marginal reduction is exact only under the added
assumption that the forward $V$'s are \emph{unpolarised} (\cref{sec:scope}); then
rotational invariance is restored and the identical one-curve-per-mass machinery
serves every parent. Because the vectors are sub-dominant to negligible for the
reach (only $J/\psi$ matters, and only for \emph{production}), this assumption
does not move the answer key; it is flagged because it is the one place the
reduction is a \emph{choice}, not a theorem.

\subsection{Two projections of one density: the rate and the shape}
\label{sec:projections}

The \emph{same} tree density \cref{eq:me} supplies the two orthogonal ingredients
the pipeline needs, as its two projections: its \emph{integral} \cref{eq:kappa}
is the \emph{rate} $\kappa_M$ (how many scalars per parent), and its
\emph{$x$-marginal} \cref{eq:factorize} is the \emph{shape} (with what energies).
In the assembled benchmark these become exactly the pinned $\kappa$ table and the
tabulated $\texttt{spectrum\_}\langle M\rangle\texttt{\_m}\langle\mphi\rangle
\texttt{.csv}$ energy spectra the meson-decay tool consumes --- which is why the
agent receives production as a \emph{provided} data product and never re-derives
the emission amplitude. (An alternative that keeps the amplitude in the loop ---
deriving the tree ME symbolically and sampling the full final state through a
generic phase-space decayer, rather than a pre-integrated curve --- is discussed
as a scope option in \cref{sec:scope,sec:scaffolding}.) The exact rates are
helicity-lifted relative to $M\to\mu\nu$: $\Gamma_3/\Gamma_2\approx0.26$--$0.30$
for the charm parents.

% =====================================================================
\section{Step 3 --- propagation: geometry, decay-in-flight, acceptance}
\label{sec:prop}

With a weighted $\phi$ flux in hand, Step~3 asks how many decay \emph{inside} the
detector and are \emph{seen}. The detector is an on-axis cylindrical decay volume
$z\in[z_{\min},z_{\max}]$ of radius $r_{\rm vol}$, followed by a two-track plane
of radius $r_{\rm det}$ at $z_{\rm det}$ (\cref{tab:geom}). Each scalar is
produced at its parent's decay vertex (below), travels in a straight line at its
lab angle, and the in-volume path segment $[L_1,L_2]$ is the general
ray--cylinder intersection \emph{from that vertex}, with
\begin{equation}
P_{\rm dec} = e^{-L_1/\lambda} - e^{-L_2/\lambda},\qquad
\lambda=\beta\gamma\,c\tau .
\end{equation}

\paragraph{Parent decay-in-flight and the production region.} The $\phi$ is not
produced at the interaction point: the parent flies from the IP and decays at
lab length $\ell\sim\mathrm{Exp}(\lambda_M)$, $\lambda_M=\beta\gamma\,c\tau_M$,
emitting the $\phi$ \emph{there}. This is $g$-independent (the parent is SM), so
it is sampled once and preserves the exact reweighting. It matters because
$c\tau_M$ spans four orders of magnitude: charm ($c\tau\!\sim\!100\,\mu$m) and
vectors ($c\tau\!\sim\!$~fm--pm) decay promptly, $\ell\!\approx\!0$; but the
\emph{kaon} has $c\tau=3.7$~m, so $\lambda_K\sim\!10$--$1000$~m --- of order the
baseline itself. The geometry then encodes a \emph{production region}
$[0,z_{\rm prod}]$: charged parents are absorbed/swept there (as behind real
shielding), so a parent counts only if it decays at $z_{\rm dec}<z_{\rm prod}$,
and $[z_{\rm prod},z_{\min}]$ is empty baseline every $\phi$ must traverse. This
is not just bookkeeping: with the absorber at the fiducial face
($z_{\rm prod}=z_{\min}$) a $\phi$ could be produced arbitrarily close to the
volume and tolerate an arbitrarily short lifetime, so the upper branch
$g_{\rm high}$ would have \emph{no well-defined value}; the gap
$z_{\min}-z_{\rm prod}$ sets a minimum survival distance and $g_{\rm high}$
converges. With $z_{\rm prod}=50$~m only the fraction $1-e^{-z_{\rm prod}
/\lambda_K}$ of kaons decay in time, which shortens the mass reach relative to
the (unphysical) prompt approximation --- the honest consequence of the kaons
being swept before most of them decay.

\paragraph{Two numerical conventions.} (1)~The decay length is sampled by
\emph{deterministic stratification} (a fixed quantile grid of the exponential per
parent, not one random draw): the upper branch is set by the sparse parents
decaying just before $z_{\rm prod}$, and stratification removes the Monte-Carlo
jitter there while keeping the exact $g$-reweighting. (2)~The visible acceptance
is evaluated at the in-volume \emph{midpoint}, a declared convention that makes
$\mathcal{A}$ \emph{$g$-independent}: the isotropic $\phi\to\mu\mu$ decay is
sampled once, both muons boosted to the lab and propagated to $z_{\rm det}$, and
both required to land within $r_{\rm det}$ moving forward (for the $\gamma\gamma$
search the two photons are collinear with the boosted $\phi$, so the requirement
is that the $\phi$ line-of-flight reach the face within $r_{\rm det}$). Computing
$\mathcal{A}$ once per mass and reweighting the $g$-scan is the $\sim\!100\times$
speed-up that makes the contour cheap. \emph{(Implementation: \texttt{parents.py}
stratified vertex sampling, \texttt{pipeline.py::segment\_from\_vertex} with the
$z_{\rm prod}$ cut; all flux assembly and propagation flow through
\texttt{propagate.py}.)}

\begin{table}[h]
\centering\footnotesize
\begin{tabular}{lcl}
\toprule
quantity & value \texttt{[PIN]} & meaning \\
\midrule
$z_{\rm prod}$ & $50$ m & production-region end (parents swept beyond) \\
$[z_{\min},z_{\max}]$ & $[112,163]$ m & decay volume \\
$r_{\rm vol}$ & $0.92$ m & volume radius \\
$z_{\rm det}$, $r_{\rm det}$ & $165$ m, $0.92$ m & two-track plane \\
$N_{\rm int}$ & $2.7\times10^{18}$ & primary $pp$ interactions \\
$N_{\rm sig}$ target & $3.0$ & zero-background contour \\
\bottomrule
\end{tabular}
\caption{Placeholder geometry/normalization. Deliberately toy (wide, close,
long) --- matching no real experiment. Levers for the ``LLP shapes''
(\cref{sec:results}): $N_{\rm int}$, $r_{\rm vol}$, $\Delta L$ set the lower
($g^4$) branch; the baseline gap $z_{\min}-z_{\rm prod}$ and the boost set the
upper branch; the open parent set and the decay-before-sweep fraction set the
mass reach.}
\label{tab:geom}
\end{table}

% =====================================================================
\section{Step 4 --- assembly, exact reweighting, and the contour}
\label{sec:assembly}\label{sec:reweight}

Step~4 multiplies the pieces together and scans the coupling. For each harvested
parent we sample the exact decay (\cref{sec:production}): $x$ from $d\Gamma/dx$ by
inverse-CDF, an isotropic direction, boost to the lab. The scalar carries a
$g^2$-\emph{stripped} weight $w_\phi=w_M\times\kappa_M(\mphi)$ --- the
per-collision parent weight \cref{eq:wnorm} times the exact branching
\cref{eq:kappa}. The full coupling dependence then factorises:
\begin{equation}
\boxed{\;N_{\rm sig}(g)
= N_{\rm int}\,g^2\sum_{i\in\text{all parents}} w_{\phi,i}\;
P_{\rm dec,i}(c\tau(g))\;\mathcal{A}_i\;}
\label{eq:reweight}
\end{equation}
One power of $g^2$ comes from production ($\kappa_M\propto g^2$, stripped and
reinstated here); the lifetime's $g^2$ enters through $c\tau(g)$ in $P_{\rm dec}$
via \cref{eq:ctau}. The kinematics ($\phi$ momenta, decay geometry) are
$g$-independent, so \emph{one} flux plus \emph{one} decay sampling covers the
entire coupling axis by reweighting --- no re-simulation per $g$. This is the
payoff of every ``$g$-independent'' convention flagged upstream (the
$g$-independent $\mathrm{BR}_{\rm vis}$ of \cref{sec:decay:total}, the
once-sampled parent vertex of \cref{sec:prop}, the midpoint acceptance).
\emph{Consistency check:} fed the same parents, this engine reproduces the
earlier analytic-kernel pipeline to the last digit ($\langle E_\phi\rangle$ and
summed weights identical, $\phi$ mass-shell to $2\times10^{-9}$).

\paragraph{The contour.} At each mass we scan $g$ on a log grid, form
$N_{\rm sig}(g)$ by \cref{eq:reweight}, and root-find $N_{\rm sig}=3$ on each
branch. The yield curve (\cref{fig:yield}) has the characteristic two-branch
shape: a $g^4$ rise at weak coupling (production $g^2\times$ linearized
$P_{\rm dec}\propto g^2$) and an exponential turnover at strong coupling (the
scalar decays before the volume). The exclusion is the closed interval
$[g_{\rm low},g_{\rm high}]$ where the peak clears $N_{\rm sig}=3$, or
\texttt{null} where it does not.

\begin{figure}[t]
\centering
\includegraphics[width=0.62\textwidth]{fig_yield_curves.pdf}
\caption{Yield curves $N_{\rm sig}(g)$ at three masses from the real flux. The
$g^4$ rise, the decay-length turnover, and the $N_{\rm sig}=3$ crossings (the two
branches) are all visible; higher masses peak lower as $K$ production shuts off.}
\label{fig:yield}
\end{figure}

% =====================================================================
\section{Results and what they teach}
\label{sec:results}

Assembling the four steps for the complete model --- both decays
(\cref{sec:decay}) and all production channels (\cref{sec:production}) --- yields
the reach, and with it the two physics lessons promised in the abstract.

\subsection{The reach: a diphoton lobe and a kaon wedge}
\label{sec:results:reach}

The honest reach at the balanced-sliver exposure $N_{\rm int}=2.7\times10^{18}$
is two disconnected pieces (\cref{fig:contour,fig:lobe}): a below-threshold
\emph{diphoton lobe} and an above-threshold \emph{kaon wedge}, meeting at the
$\sim\!10^3\times$ width cliff at $2m_\mu$ (\cref{sec:decay:total}). The
above-threshold band is
\begin{center}\footnotesize
\begin{tabular}{lccccccc}
\toprule
$\mphi$ [GeV] & 0.213 & 0.219 & 0.226 & 0.237 & 0.244 & 0.253 & 0.264 \\
\midrule
$g_{\rm low}\,[10^{-7}]$  & 7.05 & 4.60 & 4.08 & 4.00 & 4.17 & 4.70 & 6.20 \\
$g_{\rm high}\,[10^{-6}]$ & 57.2 & 13.1 & 6.91 & 3.83 & 2.86 & 2.03 & 1.25 \\
\bottomrule
\end{tabular}
\end{center}
and \texttt{null} for $\mphi\gtrsim0.27$~GeV. Both branches are smooth in mass ---
the stratified decay length (\cref{sec:prop}) removes the jitter the sparse
near-$z_{\rm prod}$ kaons would otherwise put on $g_{\rm high}$. Below threshold
the loop lifetime gives a broad lobe reaching down to $g\sim3\times10^{-6}$;
its upper edge is bounded by the baseline gap near threshold but runs to
non-perturbative $g\gtrsim1$ toward $\mphi\to0$, as $\Gamma_{\gamma\gamma}\propto
\mphi^3$ collapses and the lifetime grows without bound (an \emph{open-topped}
lobe).

\begin{figure}[t]
\centering
\includegraphics[width=\textwidth]{fig_lobe_vs_wedge.pdf}
\caption{\textbf{Left:} self-consistent contours (total width $+\,
\mathrm{BR}_{\rm vis}$): the $\gamma\gamma$ lobe below $2m_\mu$ and the $\mu\mu$
wedge above it, from the same model and geometry. \textbf{Right:} the width
partition --- $\mathrm{BR}_{\gamma\gamma}=1$ below threshold collapses as
$\mu\mu$ opens, the $\sim\!10^3\times$ jump that produces the lobe's threshold
cliff.}
\label{fig:lobe}
\end{figure}

\begin{figure}[t]
\centering
\includegraphics[width=\textwidth]{fig_leptonic_contour.pdf}
\caption{\textbf{Left:} the above-threshold reach contour in $(\mphi,g)$ from the
real Pythia flux and exact matrix element --- the two-branch sliver.
\textbf{Right:} production composition, showing the $K\to$ charm transition
across the $K$ kinematic cutoff.}
\label{fig:contour}
\end{figure}

\subsection{Why a wedge, not a lobe --- and the literature}
\label{sec:results:wedge}

The above-threshold contour is a \emph{wedge}: widest at $2m_\mu=0.2113$~GeV
(spanning $g\in[2\times10^{-7},3\times10^{-4}]$) and pinching to the high-mass
reach edge, not a closed lobe. The reason is specific to this model. Toward
threshold $\beta^3\to0$, so $\Gamma\propto g^2\mphi\beta^3\to0$ and
$c\tau\to\infty$; with $\mathrm{BR}_{\rm vis}\to1$ there is no competing decay to
cap the reach, so the long lifetime \emph{widens} the band (smaller
$g_{\rm low}$, larger $g_{\rm high}$) rather than closing it, while kaon
production stays strong (the $K$ cutoff is at $0.388$~GeV). The low-mass side is
therefore a hard \emph{kinematic wall} at $2m_\mu$, where $\phi\to\mu\mu$ closes.

This exact model is studied in Abraham \emph{et al.}, ``The LHC as a TeV Muon
Beam Dump'' (arXiv:2501.09071), with the same Yukawa $-g_S S\bar\mu\mu$ and the
same dominant production (rare meson decays; ``production from Kaons and $\omega$
mesons dominate''). Crucially they work \emph{below} threshold, $m_S<2m_\mu$,
where $S\to\gamma\gamma$ (loop, long-lived) gives a \emph{closed lobe}; above
threshold, where $\phi\to\mu\mu$ opens, they note the region is largely
constrained and leave it aside. The lobe and the wedge are thus the two regimes
of one model split at $2m_\mu$:
\begin{center}\footnotesize
\begin{tabular}{lllll}
\toprule
region & decay & width & lifetime & contour \\
\midrule
$m<2m_\mu$ (their focus) & $S\to\gamma\gamma$ (loop) & tiny & very long & closed lobe \\
$m>2m_\mu$ (\textbf{this benchmark}) & $\phi\to\mu\mu$ (tree) & large & short & threshold-bounded wedge \\
\bottomrule
\end{tabular}
\end{center}
The physics is thus fully defensible (a real, actively-studied model), while the
$m>2m_\mu$ region plus the toy geometry keeps our contour distinct from any
published result (anti-recall).

\subsection{The production staircase and the high-mass frontier}
\label{sec:results:staircase}

Adding the vector channels (\cref{sec:vector}) reveals the second lesson. The
dominant \emph{producer} steps cleanly with mass (\cref{fig:staircase} right,
\cref{fig:composition}): $K$ ($\mphi\lesssim0.39$) $\to$ $D_s$ (leptonic,
$0.4$--$1.4$) $\to$ $J/\psi$ ($1.4$--$2.4$) $\to$ $\psi(2S)$, a real
\emph{production} staircase, with the clean, exactly-computable $J/\psi$ channel
at its top (from $\mathrm{BR}(J/\psi\to\mu\mu)=5.9\%$ alone). The light vectors
$\rho/\omega/\phi$, though as abundant as kaons, are $\alpha^2$-suppressed by
$\mathrm{BR}(V\to\mu\mu)\sim10^{-4}$ and sit $>\!100\times$ below the leading
channel everywhere --- omitting them is a sub-per-mille effect.

But the \emph{reach} tells a soberer story (\cref{fig:staircase} left): at this
exposure only the kaon wedge clears $N_{\rm sig}=3$. The center panel plots the
\emph{crest} of $N_{\rm sig}(g)$ versus mass; it crosses $3$ at
$\mphi\approx0.32$, already below the kaon cutoff, and never returns. The
mechanism is instructive: it is \emph{not} that heavier $\phi$ decay too soon
(at the crest, $g$ is tuned so $\lambda$ lands in the volume,
$P_{\rm dec}\!\approx\!0.14$, roughly mass-independent). The crest sinks because
(i) $\kappa_K$ \emph{collapses} as $\mphi$ nears the kaon endpoint --- the $\phi$
takes almost all the energy and the $\mu\nu$ phase space vanishes --- and above
$0.388$ kaons cannot make $\phi$ at all; and (ii) the two-track acceptance falls,
since a heavier, less-boosted $\phi$ opens its $\mu^+\mu^-$ wider. Charm and
$J/\psi$, $\sim\!100\times$ rarer, never lift the crest back over $3$.

\begin{figure}[t]
\centering
\includegraphics[width=\textwidth]{fig_staircase.pdf}
\caption{\textbf{Left:} all-channel $\mu\mu$ reach at $N_{\rm int}=2.7\times
10^{18}$ --- only the kaon wedge clears $N_{\rm sig}=3$. \textbf{Center:} the
peak yield (optimized over $g$ at each mass) crests at threshold and slides under
$N_{\rm sig}=3$ at $\sim\!0.32$~GeV, never recovering --- why the band closes.
\textbf{Right:} production composition --- the dominant producer steps
$K\to D_s\to J/\psi\to\psi(2S)$ (the staircase is real in \emph{production}).}
\label{fig:staircase}
\end{figure}

\begin{figure}[t]
\centering
\includegraphics[width=0.72\textwidth]{fig_production_composition.pdf}
\caption{$g^2$-stripped $\phi$ production per collision, rate$\times\kappa$, for
all eight channels (pseudoscalar solid, vector dashed). The light vectors sit
$>\!100\times$ below the leading channel everywhere; only $J/\psi$ (thick) rises
to dominate above $\sim\!1.4$~GeV.}
\label{fig:composition}
\end{figure}

\subsection{The geometry scan: one pipeline, many LLP shapes}
\label{sec:results:geomscan}

Because the forward acceptance is folded downstream (\cref{sec:pythia}), the
\emph{same} flux and toolchain produce different contour morphologies under
different pinned geometry/normalization --- the reinterpretation thesis in
action. Five configurations, now with \emph{both} channels (each a
$\gamma\gamma$ lobe below $2m_\mu$ plus a $\mu\mu$ wedge above):
\begin{center}\scriptsize
\begin{tabular}{lccccccc}
\toprule
regime & $z_{\rm prod}$ & $z_{\min}$ & $x_{\rm off}$ & $N_{\rm int}$ & $\gamma\gamma$ $g_{\rm lo}$ & $\mu\mu$ wedge $[g_{\rm lo},g_{\rm hi}]$ & reach \\
\midrule
A: balanced sliver          &  50 & 112 & 0 & $2.7{\times}10^{18}$ & $3.5{\times}10^{-6}$ & $[4.0{\times}10^{-7},\,5.7{\times}10^{-5}]$ & $0.26$ \\
B: deep tower (hi-lumi)     &  50 & 112 & 0 & $1.0{\times}10^{20}$ & $1.4{\times}10^{-6}$ & $[1.3{\times}10^{-7},\,8.8{\times}10^{-5}]$ & $0.40$ \\
C: near short-baseline      &  15 &  32 & 0 & $2.7{\times}10^{18}$ & $3.4{\times}10^{-6}$ & $[3.7{\times}10^{-7},\,1.3{\times}10^{-4}]$ & $0.30$ \\
D: far compact (FASER-like) &  50 & 480 & 0 & $3.0{\times}10^{18}$ & $6.0{\times}10^{-6}$ & $[1.5{\times}10^{-6},\,1.6{\times}10^{-5}]$ & $0.213$ \\
E: off-axis (A, $1$ m)      &  50 & 112 & 1 & $2.7{\times}10^{18}$ & $3.7{\times}10^{-6}$ & $[4.7{\times}10^{-7},\,4.5{\times}10^{-5}]$ & $0.245$ \\
\bottomrule
\end{tabular}
\end{center}
The knobs behave as \cref{sec:prop} predicts: the lower branch scales roughly
$g_{\rm low}\propto(N_{\rm int}\,\Delta L)^{-1/4}$ (the hi-lumi tower B plunges
deepest); the $\mu\mu$ upper branch is set by the \emph{baseline gap}
$z_{\min}-z_{\rm prod}$, not $z_{\min}$ alone (near detector C reaches the largest
coupling, far D the smallest); and only the hi-lumi B pushes the wedge past the
kaon region to $\sim\!0.40$~GeV. These are the marquee ``LLP shapes'': the
balanced banana (A), the tall deep-reach tower (B), the elevated near-baseline
(C), the narrow far sliver (D), the softened off-axis (E) --- each a lobe$+$wedge
pair, each a pinned-data-product swap through the identical pipeline.

\begin{figure}[t]
\centering
\includegraphics[width=\textwidth]{fig_geometry_scan.pdf}
\caption{Geometry scan, both channels: one flux and one toolchain yield five
qualitatively distinct morphologies under different pinned geometry/normalization,
each a $\gamma\gamma$ lobe plus a $\mu\mu$ wedge. Lower-branch depth is set by
$N_{\rm int}$ and the volume, the $\mu\mu$ upper branch by the baseline gap
$z_{\min}-z_{\rm prod}$; the hi-lumi tower (B) extends the wedge past the kaon
region, and the off-axis E is a softened A.}
\label{fig:geomscan}
\end{figure}

\subsection{The load-bearing finding}
\label{sec:finding}

Both lessons combine into one sentence: \textbf{at LHC-forward exposure the reach
is limited by absolute forward rate, not by channel availability.} The dominant
producer does climb $K\to D_s\to J/\psi$ with mass, and the channel at the top of
that staircase, $J/\psi\to\mu\mu\phi$, is \emph{exactly computable} --- the clean
high-mass producer one might have hoped semileptonic charm would be, without its
form-factor baggage. But neither leptonic charm nor $J/\psi$ clears
$N_{\rm sig}=3$ here: $D_s$ falls $\sim\!20\times$ short near $0.4$~GeV, $J/\psi$
far more above $1.5$~GeV. Converting the production staircase into a \emph{reach}
staircase needs beam-dump-scale exposure (the regime-B lever), and even then the
accessible extension is the charm window $0.4$--$1.8$~GeV, lifted into reach by
the \emph{semileptonic} charm channel $D\to X\mu\nu\phi$ (muon flux
$\sim\!500\times$ the leptonic branching); the $J/\psi$ window above the $D_s$
cutoff stays out of reach even then --- a production frontier, not a reach one.
Semileptonic charm is the remaining brick, deferred (\cref{sec:scope}).

% =====================================================================
\section{Scope, omissions, and toy approximations}
\label{sec:scope}

The answer key aims to be \emph{correct and well-defined}, not \emph{complete}:
it is a pinned toy, so what matters is that every included piece is right and
every excluded piece is declared. This section is the authoritative list.

\paragraph{Included.} \emph{Decay}: $\phi\to\mu\mu$ (tree) and $\phi\to
\gamma\gamma$ (one-loop), the exact total width. \emph{Production}: the exact
$M\to\mu\nu\phi$ for $\pi,K,D,D_s$ and the exact $V\to\mu\mu\phi$ for
$\rho^0,\omega,\phi(1020),J/\psi,\psi(2S)$.

\paragraph{Genuine omissions.}
\begin{itemize}
\item \textbf{Muon bremsstrahlung} $\mu N\to\mu N\phi$ --- \emph{the leading
omission}. Real forward muons (copious, from $\pi/K\to\mu\nu$) radiate $\phi$ in
material; a free on-shell muon cannot, so this needs the nuclear field. It is the
reference paper's whole premise (Abraham \emph{et al.}) and \emph{the} production
mechanism for a beam-dump sibling. We omit it: our setting is collider-forward
(near-vacuum), so it is a correction here. It is also material-transport physics,
which does not share the exactly-re-derivable property of the vacuum pipeline.
\item \textbf{Semileptonic charm} $D\to X\mu\nu\phi$ --- deferred (hadronic form
factors) and null-by-rate at our exposure, but the honest route to the
charm-window reach at beam-dump luminosity.
\item \textbf{$\eta,\eta'\to\mu^+\mu^-\phi$} and other light-meson dimuon Dalitz
channels --- $\alpha^2$-suppressed, negligible, exactly like the light vectors.
\end{itemize}

\paragraph{Toy approximations (correct-as-modelled, deliberately simple).}
Forward charm and prompt-quarkonium fluxes carry real Pythia modelling
uncertainty --- absorbed because our run \emph{is} the reference. The statistical
rule $N_{\rm sig}=3$, zero background is apt for the displaced $\mu\mu$ vertex but
\emph{not} for $\gamma\gamma$ (which has real backgrounds), so the lobe's
absolute reach would tighten under a proper significance rule. The detector is
idealised: no muon energy threshold, no efficiency/resolution, straight-line
propagation (no $B$-field). Vector parents are taken unpolarised (so the $\phi$
direction is isotropic, \cref{sec:reduction}). The $\gamma\gamma$ acceptance is
the collinear-photon placeholder.

\paragraph{Production as a spectrum vs.\ as a live decay (a scope option).} The
pipeline samples production from a \emph{pre-computed} energy spectrum
$d\Gamma/dx$ (\cref{eq:factorize}): the emission matrix element is integrated
once, offline, into one curve per (parent, mass); the tool then only draws $x$,
throws an isotropic direction, and boosts. This is exact for the LLP flux and
fast, but it (a)~bakes in the isotropy assumption (exact for spin-0 parents, an
approximation for polarised vectors); (b)~needs a tabulated curve per (parent,
mass); and (c)~exposes only the $\phi$, discarding the recoil $\mu\nu$. A more
general alternative keeps the amplitude \emph{in the loop}: derive the tree
production ME symbolically (the single-diagram $M\to\mu\nu\phi$ and two-diagram
$V\to\mu\mu\phi$ are $1\to3$ \emph{tree} processes, within scope of the exact
symbolic tools, unlike the $\phi\to\gamma\gamma$ loop) and feed it to a generic
phase-space decayer that boosts the parent to rest, samples the full final state
over phase space weighted by $|\Ma|^2$, and boosts every product back to the lab.
That route is more general (any $1\to N$, any mass, full final state, correct
vector polarisation) and makes production \emph{derived-and-verifiable} rather
than a provided table --- at the cost of a per-event ME evaluation and of
\emph{widening the task scope} to include the production derivation the current
key deliberately provides. The trade is developed in \cref{sec:scaffolding}.

% =====================================================================
\section{From solution to scaffolding: what tools this class of problem wants}
\label{sec:scaffolding}

Reading the solution back as an interface map, each pipeline step asks a
different thing of a toolchain --- which is the point of laying the note out
step-wise. This class of problem (a BSM LLP produced in meson decays, propagated
through a forward volume) decomposes into:

\begin{center}\footnotesize
\begin{tabular}{lll}
\toprule
step & what it computes & natural tool / existing analog \\
\midrule
theory input (\cref{sec:decay}) & tree \& loop widths, $\mathrm{BR}$, $c\tau$ & symbolic amplitude tools;
  the loop needs by-hand input \\
Step 1 flux (\cref{sec:pythia}) & forward parent spectra & event generator run-card driver (Pythia) \\
Step 2 production (\cref{sec:production}) & parent $\to\phi$ ME, $\kappa$, spectrum & \emph{either} a provided-spectrum
  sampler \\
 & & \emph{or} ME derivation $+$ phase-space decayer \\
Step 3 propagation (\cref{sec:prop}) & decay-in-flight, $P_{\rm dec}$, $\mathcal{A}$ & geometry / decay-in-volume tool \\
Step 4 assembly (\cref{sec:assembly}) & $g^2$ reweight, root-find contour & thin numerical driver (no special tool) \\
\bottomrule
\end{tabular}
\end{center}

Two boundaries carry the design tension. \emph{The theory rung} (\cref{sec:decay})
splits cleanly: the tree $\mu\mu$ width is a $1\to2$ symbolic-tool target, but the
one-loop $\gamma\gamma$ width is outside a tree-only symbolic engine --- so a
faithful benchmark either supplies it or accepts a by-hand derivation, and should
not pretend a tree tool produces it. \emph{The production step}
(\cref{sec:production}) is the spectrum-vs-decayer fork of \cref{sec:scope}: the
provided-spectrum route mirrors the tabulated forward-LLP tools (FORESEE,
ALPINIST) and is fast and reinterpretation-friendly; the ME-derivation route
mirrors the general generators (a custom-model \textsc{EvtGen} for meson decays;
FeynRules/UFO~$\to$~MadGraph for parton-level production) and is more general and
more verifiable. Notably the second route's scaffolding is largely pre-existing
(symbolic-derivation, FeynRules, and MadGraph wrappers already live in the
toolset), so the choice is one of \emph{task scope} --- provide production or
require its derivation --- more than of missing machinery. This map, not any
single number, is the reusable output of the solution.

% =====================================================================
\section{What is pinned, what is computed, and reproduction}
\label{sec:repro}

\textbf{Computed from first principles:} the tree width \cref{eq:width} and the
one-loop $\phi\to\gamma\gamma$ width \cref{eq:gammagg}, the exact $M\to\mu\nu\phi$
and $V\to\mu\mu\phi$ matrix elements and their $\kappa$
(\cref{sec:me,sec:vector}), the forward flux (Pythia, \cref{sec:pythia}), and
everything downstream. \textbf{Pinned inputs \texttt{[PIN]} (synthetic,
anti-recall):} $\sqrt{s}$, $\sigma_{\rm inel}$, the geometry and $N_{\rm int}$
(\cref{tab:geom}), the mass grid, the statistical rule, and the
$\mathrm{BR}(V\to\mu\mu)$ set.

Reproduce with (hepbench conda env):
\begin{lstlisting}
# 1. forward all-channel flux (K/D/Ds + rho/omega/phi/Jpsi/psi2S; ~14 min,
#    ~300 MB for 1e6 events)
python _shared/generate_parent_flux.py --n-soft 1000000 --n-hard 0 \
       --theta-max 0.05 --seed 2026 --out parent_flux_allchan_1M.npz
# 2. leptonic wedge, geometry scan, flux spectra
python _shared/notes/make_leptonic_figs.py parent_flux_allchan_1M.npz
# 3. diphoton lobe vs mumu wedge (self-consistent total width + BR_vis)
python _shared/notes/make_diphoton_figs.py parent_flux_allchan_1M.npz
# 4. all-channel reach + peak-N_sig + production staircase (K -> Ds -> J/psi)
python _shared/notes/make_staircase_fig.py parent_flux_allchan_1M.npz
# 5. geometry scan, both channels (lobe + wedge x 4 regimes)
python _shared/notes/make_geomscan_fig.py parent_flux_allchan_1M.npz
\end{lstlisting}
Physics validation: \texttt{\_shared/production.py::validate} (trace vs.\ spinor,
quadrature vs.\ Dalitz, frame invariance, $\phi\to\mu\mu$ closed form vs.\
numerics, and the diphoton closed form vs.\ its Feynman-parameter integral) and
\texttt{\_shared/production\_vector.py::validate} (Ward identity, 2-body
normalization). This note documents the state \emph{before} wiring into the
benchmark \texttt{ground\_truth/derivation.py} and \emph{before} the semileptonic
brick; both follow once the physics here is signed off.

\end{document}
